% !TeX root = ../../report.tex
\subsection{Evolutionary Algorithms}\label{sec:ea}

Evolutionary algorithms (EAs) are a class of stochastic search and metaheuristic optimisation techniques inspired by biological evolution and selection. Unlike traditional algorithms that iterate on a single solution, EAs maintain a population of individuals, where each individual represents a potential solution to a problem \citep{yedjour2024}. The population evolves over generations using three primary operators: selection, crossover, and mutation.

Rule extraction is a post-hoc explainability technique used to identify the underlying decision-making logic of opaque black-box models such as ANNs \citep{yedjour2024}. It is fundamentally a multi-objective optimisation problem that must balance fidelity to the black-box model, predictive accuracy on unseen data, and rule comprehensibility \citep{yedjour2024}. Deterministic rule-extraction methods such as decision trees impose fixed inductive biases on the search (e.g., axis-aligned splits, depth constraints, and greedy coverage strategies), which can limit their ability to navigate the combinatorial complexity of high-dimensional genomic feature spaces or optimise competing objectives simultaneously \citep{arrieta2020,yedjour2024}. EAs are well suited to this problem because they maintain a diverse population of candidate rule sets, can be configured with flexible multi-objective fitness functions, and can escape local optima through crossover and mutation \citep{yedjour2024}. Among EA variants, genetic algorithms (GA) and non-dominated sorting genetic algorithm II (NSGA-II) are selected for this project: GA provides a single-objective baseline for rule-extraction quality, while NSGA-II explicitly evolves a Pareto front of non-dominated rule sets, enabling structured trade-off analysis between fidelity and comprehensibility \citep{deb2002}.